\documentclass[10pt, letterpaper]{article}

% Packages:
\usepackage[
    ignoreheadfoot, % set margins without considering header and footer
    top=1.6 cm, % seperation between body and page edge from the top
    bottom=1.6 cm, % seperation between body and page edge from the bottom
    left=1.75 cm, % seperation between body and page edge from the left
    right=1.75 cm, % seperation between body and page edge from the right
    footskip=1.0 cm, % seperation between body and footer
    % showframe % for debugging 
]{geometry} % for adjusting page geometry
\usepackage{titlesec} % for customizing section titles
\usepackage{tabularx} % for making tables with fixed width columns
\usepackage{array} % tabularx requires this
\usepackage[dvipsnames]{xcolor} % for coloring text
\definecolor{primaryColor}{RGB}{0, 0, 0} % define primary color
\usepackage{enumitem} % for customizing lists
\usepackage{fontawesome5} % for using icons
\usepackage{amsmath} % for math
\usepackage[
    pdftitle={Lei Gong - Tailored Resume},
    pdfauthor={Lei Gong},
    pdfcreator={LaTeX with RenderCV},
    colorlinks=true,
    urlcolor=primaryColor
]{hyperref} % for links, metadata and bookmarks
\usepackage[pscoord]{eso-pic} % for floating text on the page
\usepackage{calc} % for calculating lengths
\usepackage{bookmark} % for bookmarks
\usepackage{lastpage} % for getting the total number of pages
\usepackage{changepage} % for one column entries (adjustwidth environment)
\usepackage{paracol} % for two and three column entries
\usepackage{ifthen} % for conditional statements
\usepackage{needspace} % for avoiding page brake right after the section title
\usepackage{iftex} % check if engine is pdflatex, xetex or luatex

% Ensure that generate pdf is machine readable/ATS parsable:
\ifPDFTeX
    \input{glyphtounicode}
    \pdfgentounicode=1
    \usepackage[T1]{fontenc}
    \usepackage[utf8]{inputenc}
    \usepackage{lmodern}
\fi

\usepackage{charter}

\usepackage{graphicx}
\usepackage{hyperref}

% Some settings:
\raggedright
\AtBeginEnvironment{adjustwidth}{\partopsep0pt} % remove space before adjustwidth environment
\pagestyle{empty} % no header or footer
\setcounter{secnumdepth}{0} % no section numbering
\setlength{\parindent}{0pt} % no indentation
\setlength{\topskip}{0pt} % no top skip
\setlength{\columnsep}{0.15cm} % set column seperation
\pagenumbering{gobble} % no page numbering

\titleformat{\section}{\needspace{4\baselineskip}\bfseries\large}{}{0pt}{}[\vspace{1pt}\titlerule]

\titlespacing{\section}{
    % left space:
    -1pt
}{
    % top space:
    0.16 cm
}{
    % bottom space:
    0.08 cm
} % section title spacing

\renewcommand\labelitemi{$\vcenter{\hbox{\small$\bullet$}}$} % custom bullet points
\newenvironment{highlights}{
    \begin{itemize}[
        topsep=0.03 cm,
        parsep=0.03 cm,
        partopsep=0pt,
        itemsep=0pt,
        leftmargin=0 cm + 10pt
    ]
}{
    \end{itemize}
} % new environment for highlights


\newenvironment{highlightsforbulletentries}{
    \begin{itemize}[
        topsep=0.10 cm,
        parsep=0.10 cm,
        partopsep=0pt,
        itemsep=0pt,
        leftmargin=10pt
    ]
}{
    \end{itemize}
} % new environment for highlights for bullet entries

\newenvironment{onecolentry}{
    \begin{adjustwidth}{
        0 cm + 0.00001 cm
    }{
        0 cm + 0.00001 cm
    }
}{
    \end{adjustwidth}
} % new environment for one column entries

\newenvironment{twocolentry}[2][]{
    \onecolentry
    \def\secondColumn{#2}
    \setcolumnwidth{\fill, 4.5 cm}
    \begin{paracol}{2}
}{
    \switchcolumn \raggedleft \secondColumn
    \end{paracol}
    \endonecolentry
} % new environment for two column entries

\newenvironment{threecolentry}[3][]{
    \onecolentry
    \def\thirdColumn{#3}
    \setcolumnwidth{, \fill, 4.5 cm}
    \begin{paracol}{3}
    {\raggedright #2} \switchcolumn
}{
    \switchcolumn \raggedleft \thirdColumn
    \end{paracol}
    \endonecolentry
} % new environment for three column entries

\newenvironment{header}{
    \setlength{\topsep}{0pt}\par\kern\topsep\centering\linespread{1.5}
}{
    \par\kern\topsep
} % new environment for the header

\newcommand{\placelastupdatedtext}{% \placetextbox{<horizontal pos>}{<vertical pos>}{<stuff>}
  \AddToShipoutPictureFG*{% Add <stuff> to current page foreground
    \put(
        \LenToUnit{\paperwidth-2 cm-0 cm+0.05cm},
        \LenToUnit{\paperheight-1.0 cm}
    ){\vtop{{\null}\makebox[0pt][c]{
        \small\color{gray}\textit{Last updated in September 2024}\hspace{\widthof{Last updated in September 2024}}
    }}}%
  }%
}%

% save the original href command in a new command:
\let\hrefWithoutArrow\href

% new command for external links:

% \immediate\write18{curl -L -o medrag-stars.png "https://img.shields.io/github/stars/Teddy-XiongGZ/MedRAG?style=social"}



\begin{document}
\newcommand{\AND}{\unskip\cleaders\copy\ANDbox\hskip\wd\ANDbox\ignorespaces}
\newsavebox\ANDbox
\sbox\ANDbox{$|$}
\begin{header}
\fontsize{23 pt}{23 pt}\selectfont Lei Gong

\vspace{3 pt}
\normalsize
\mbox{Charlottesville, VA} \kern 5.0 pt \AND \kern 5.0 pt
\mbox{\hrefWithoutArrow{mailto:nkw3mr@virginia.edu}{nkw3mr@virginia.edu}} \kern 5.0 pt \AND \kern 5.0 pt
\mbox{\hrefWithoutArrow{tel:+17657719032}{(+1) 765-771-9032}} \kern 5.0 pt \AND \kern 5.0 pt
\mbox{\hrefWithoutArrow{https://github.com/LeiGong0125Carrot}{GitHub}} \kern 5.0 pt \AND \kern 5.0 pt
\mbox{\hrefWithoutArrow{https://scholar.google.com/citations?user=8QLAkKMAAAAJ}{Google Scholar}}
\end{header}

\vspace{0.06 cm}
\section{SUMMARY}
\begin{onecolentry}
NLP and multimodal ML researcher working on retrieval-augmented information extraction, language-model adaptation, and rigorous experimental evaluation.
\end{onecolentry}
\section{EDUCATION}
\begin{onecolentry}
\textbf{University of Virginia} \hfill Charlottesville, VA \\
M.S. Computer Science, GPA 3.91/4.00 \hfill Expected December 2026
\end{onecolentry}
\vspace{0.04 cm}
\begin{onecolentry}
\textbf{Purdue University} \hfill West Lafayette, IN \\
B.S. Computer Science; Second Major in Applied Statistics \hfill May 2023
\end{onecolentry}
\section{TECHNICAL SKILLS}
\begin{onecolentry}
\textbf{ML and NLP:} Python, PyTorch, transformers, retrieval-augmented learning, language models, information extraction \\ \textbf{Data and Knowledge:} SQL, Neo4j, knowledge graph construction, clinical NLP, MetaMap
\end{onecolentry}
\section{RESEARCH EXPERIENCE}
\begin{onecolentry}
\textbf{Graduate Research Assistant, University of Virginia} \hfill Jan 2024 -- Present
\begin{highlights}
\item Conducted research in retrieval-augmented NLP, semantic enrichment, and language-model evaluation, developing and evaluating machine learning methods with clinical notes, structured EHR features, and external biomedical knowledge.
\item Built and experimentally evaluated research models for healthcare prediction and social determinants of health extraction; collaborated with research coauthors on model development and publication.
\end{highlights}
\end{onecolentry}
\section{SELECTED PROJECTS}

\begin{twocolentry}{IEEE Big Data 2024}
\textbf{Retrieval-Augmented SDoH Extraction}
\end{twocolentry}
\begin{onecolentry}
\begin{highlights}
\item Developed a retrieval-augmented clinical NLP framework integrating biomedical literature with EHR narratives using cross-source attention and gated feature fusion.
\item Compared against six baselines on a MIMIC-IV SDoH benchmark; improved macro-F1 by 2.91 percentage points and macro-AUC by 2.49 percentage points over Hierarchical RoBERTa.
\end{highlights}
\end{onecolentry}
\vspace{0.06 cm}

\begin{twocolentry}{arXiv 2026}
\textbf{OPD-Aha: Multimodal On-Policy Distillation}
\end{twocolentry}
\begin{onecolentry}
\begin{highlights}
\item Collaborated on multimodal on-policy distillation research, contributing to algorithm development discussions and Qwen3.5-4B/9B post-training with visual preference-based target reconstruction.
\item Contributed to experimental validation and benchmark evaluation; the OPD-Aha framework achieved +4.0/+3.2 percentage-point accuracy gains over Vision-OPD across six visual reasoning benchmarks.
\end{highlights}
\end{onecolentry}
\vspace{0.06 cm}

\begin{twocolentry}{AMIA 2024}
\textbf{Knowledge-Augmented LLM for SDoH Extraction}
\end{twocolentry}
\begin{onecolentry}
\begin{highlights}
\item Integrated UMLS biomedical semantic types into clinical narratives to improve language-model-based social determinants of health extraction under limited data.
\item Evaluated FLAN-T5 models at multiple scales on MIMIC-III SDoH data, increasing macro-F1 from 0.8423 to 0.8824.
\end{highlights}
\end{onecolentry}
\vspace{0.06 cm}

\begin{twocolentry}{2026 -- Present}
\textbf{SCoPe: Semantic Context-Adaptive Prototypical Learning}
\end{twocolentry}
\begin{onecolentry}
\begin{highlights}
\item Developed a context-adaptive representation learning framework with semantic unit aggregation, class-conditional prototypical learning, and KL regularization for interpretable clinical outcome prediction.
\item Evaluated on MIMIC-III, MIMIC-IV, and Hyperpartisan benchmarks; improved mortality prediction F1 from 0.268 to 0.308 over Clinical-ModernBERT.
\end{highlights}
\end{onecolentry}
\vspace{0.06 cm}
\section{SELECTED PUBLICATIONS}
\begin{onecolentry}
Coauthor, \textit{Reward Is Enough: LLMs Are In-Context Reinforcement Learners} (\textbf{ICLR 2026}); coauthor, \textit{OPD-Aha} (arXiv 2026).
\end{onecolentry}
\end{document}
